\begin{corollary}[5.2.3]
\EGAReaderTarget{ega4.c04.cor.5-2-3}{IV.2.5.2.3-fr}
Sean $X$ un preesquema localmente de tipo finito sobre un cuerpo $k$ y $x$ un
punto de $X$. Se tiene
\begin{equation}
\dim_x(X)=\dim(\mathcal O_x)+\operatorname{deg.tr}_k \boldsymbol{k}(x).
\tag{5.2.3.1}\label{IV.2.5.2.3.1-fr}
\end{equation}
\end{corollary}

\begin{proof}
Se sabe que existe una vecindad abierta $U$ de $x$ en $X$ tal que
$\dim_x(X)=\dim(U)$ $(\mathbf{0},14.1.4.1)$, y se puede suponer que las componentes
irreducibles de $U$ sean los $X_i\cap U$, donde los $X_i$ son las
componentes irreducibles de $X$ que contienen a $x$; como $U\cap X_i$

\oldpage[IV-2]{91}

es denso en $X_i$, de (4.1.1.3) se sigue que
$\dim(X_i)=\dim(U\cap X_i)$, por lo que se tiene
$\dim_x(X)=\sup_i(\dim(X_i))$. Además, los ideales primos minimales de
$\mathcal O_x$ corresponden a los puntos genéricos de los $X_i$, por lo que,
por $(\mathbf{0},16.1.1.1)$, se tiene
$\dim(\mathcal O_{X,x})=\sup_i(\dim(\mathcal O_{X_i,x}))$. La cuestión se reduce así
al caso en que $X$ es \emph{irreducible}; como $X$ es biequidimensional por
(5.2.1), se tiene
$\dim(X)=\dim(\overline{\{x\}})+\operatorname{codim}(\overline{\{x\}},X)$
$(\mathbf{0},14.3.5.1)$, y se sabe que
$\dim(\overline{\{x\}})=\operatorname{deg.tr}_k \boldsymbol{k}(x)$ por (5.2.1) y
$\operatorname{codim}(\overline{\{x\}},X)=\dim(\mathcal O_x)$ por (5.1.2).
\end{proof}

\begin{corollary}[5.2.4]
\label{IV.2.5.2.4-fr}
Sean $X$ un preesquema localmente de tipo finito sobre un cuerpo $k$,
$\mathcal L$ un $\mathcal O_X$-módulo invertible y $f$ una sección de
$\mathcal L$ sobre $X$, tal que el conjunto $Y$ de los $x\in X$ para los
cuales $f(x)=0$ sea raro en $X$. Entonces se tiene
$\dim(Y)\leq\dim(X)-1$; los dos miembros de esta desigualdad son iguales si
$Y$ corta toda componente irreducible de dimensión máxima de $X$.
\end{corollary}

\begin{proof}
Si $(X_\alpha)$ es la familia de las componentes irreducibles de $X$, se
tiene
\[
\dim(Y)=\sup_\alpha(\dim(Y\cap X_\alpha))
\]
y, por tanto, la cuestión se reduce al caso en que $X$ es irreducible (pues
$Y\cap X_\alpha$ es raro en $X_\alpha$, ya que todo $X_\alpha$ tiene interior
no vacío en $X$). Basta considerar el caso en que $Y\ne\varnothing$; entonces,
para todo punto maximal $x$ de $Y$, se tiene (puesto que $X$ es
biequidimensional)
$\dim(\overline{\{x\}})=\dim(X)-\operatorname{codim}(\overline{\{x\}},X)$,
y, puesto que $Y$ es raro en $X$, se tiene
\[
\operatorname{codim}(\overline{\{x\}},X)=1
\] por (5.1.8); de aquí el
corolario.
\end{proof}

\begin{remarks}[5.2.5]
\label{IV.2.5.2.5-fr}
\begin{enumerate}
\item[(i)] A diferencia de lo que ocurre con los $k$-preesquemas algebraicos,
un preesquema localmente noetheriano $X$ no es necesariamente catenario
(cf. (5.6.11)); sin embargo, lo es si todos sus anillos locales $\mathcal O_x$
son cocientes de anillos regulares $(\mathbf{0},16.5.12)$ y, en particular, si $X$ es
\emph{regular} (\textbf I, 4.1.4). No obstante, ni siquiera un esquema (íntegro)
$X=\operatorname{Spec}(A)$, donde $A$ es un anillo \emph{regular}, es
necesariamente \emph{biequidimensional}; dicho de otro modo $(\mathbf{0},14.3.3)$,
las codimensiones en $X$ de los distintos puntos cerrados de $X$ no son
necesariamente las mismas. Por ejemplo, sean $B$ un anillo de valoración
discreta, $\mathfrak m=B\pi$ su ideal maximal, $k=B/\mathfrak m$ su cuerpo
residual y $K$ el cuerpo de fracciones de $B$; sea $A$ el anillo de polinomios
$B[T]$. En $A$ hay ideales maximales \emph{de altura} $2$, por ejemplo
$(\pi)+(T)$; pero también hay ideales maximales \emph{de altura} $1$, por
ejemplo el ideal principal $(\pi T-1)$: en efecto, un ideal primo principal
$\ne0$ es de altura $1$ (5.1.8); por otra parte,
$A/(\pi T-1)$ es isomorfo al anillo de fracciones $B_\pi$
$(\mathbf{0}_{\mathrm I},1.2.3)$, que aquí no es otro que $K$, lo que prueba que el
ideal $(\pi T-1)$ es maximal y de altura $1$.

\item[(ii)] Cuando $X$ es un preesquema localmente de tipo finito sobre un
cuerpo, se ha visto (4.1.1.3) que $\dim(U)=\dim(X)$ para todo abierto $U$
denso en todas partes en $X$. Este resultado \emph{no se extiende} a los
preesquemas noetherianos regulares, aunque sean biequidimensionales. Por
ejemplo, sean $B$ un anillo de valoración discreta y $\mathfrak m$ su ideal
maximal; $X=\operatorname{Spec}(A)$ tiene dos puntos, el ideal $(0)$ y
$\mathfrak m$, siendo este último el único punto cerrado; se tiene
$\dim(X)=1$, pero el conjunto abierto $U=\{(0)\}$ es de dimensión $0$
(cf. \S~10).
\end{enumerate}
\end{remarks}
